% ============================================================
%  Experiments section — NeurIPS style
%  % ============================================================
%  Experiments section — NeurIPS style
%  % ============================================================
%  Experiments section — NeurIPS style
%  % ============================================================
%  Experiments section — NeurIPS style
%  \input{experiments} from your main .tex file
% ============================================================

\section{Experiments}

\subsection{Dataset and Task}

We evaluate on \textbf{EgoGazeVQA}~\cite{egogaze}, a dataset of egocentric video clips from
Ego4D~\cite{ego4d} paired with gaze annotations and multiple-choice questions (MCQ).
Each sample consists of a short video clip, a five-way MCQ, and per-frame gaze coordinates
recorded via eye-tracking.
Questions are categorised into three reasoning types: \emph{causal} (why did the person do X),
\emph{spatial} (where is object Y), and \emph{temporal} (what happened before/after Z).

We sample \textbf{30 clips} with a balanced mix of question types (seed 42)
and extract \textbf{4 frames} per clip at uniform temporal intervals.
Ground-truth scene descriptions are obtained by concatenating the Ego4D
narration annotations~\cite{ego4d} for each clip's frame range.

\subsection{Model and Implementation Details}

We use \textbf{Qwen2-VL-2B-Instruct}~\cite{qwen2vl} as our base vision-language model,
loaded in \texttt{bfloat16} on a single NVIDIA A100-SXM4-40GB GPU.
We modify the model's transformer decoder to support \emph{visual token pruning}:
at a designated layer $\ell$, a fraction $(1 - r)$ of visual tokens are dropped
from the key-value cache before processing the remaining decoder layers,
reducing the effective sequence length for all subsequent layers.

We evaluate four pruning strategies:
\begin{itemize}
  \item \textbf{No pruning} (baseline): all visual tokens retained throughout.
  \item \textbf{Random}: tokens selected uniformly at random.
  \item \textbf{Gaze-guided}: tokens scored by a 2-D Gaussian centred on the
        per-frame gaze coordinate; top-$r$ fraction retained.
  \item \textbf{Text-guided}: tokens scored by rater-token attention
        (SparseVLM-style~\cite{sparsevlm}); top-$r$ fraction retained.
  \item \textbf{Combined}: convex combination of gaze and text scores,
        $s = \alpha \cdot s_{\text{text}} + (1-\alpha) \cdot s_{\text{gaze}}$.
\end{itemize}

Unless otherwise stated, we fix the pruning layer to $\ell = 10$
(selected via a prior layer ablation) and the keep ratio to $r = 0.5$.
The model is prompted to reason step-by-step before outputting
\texttt{Answer: X}, allowing us to capture the full chain-of-thought (CoT).

\subsection{Evaluation Metrics}

We report the following metrics across all conditions:

\paragraph{MCQ Accuracy.}
Exact-match accuracy between the model's predicted answer letter and the
ground-truth option, computed over all 30 samples per condition.
We break accuracy down by question type (causal / spatial / temporal).

\paragraph{Decode Throughput.}
Mean decode latency in milliseconds per generated token (ms/tok),
averaged over all samples. This measures the inference speedup from
reduced KV-cache size.

\paragraph{Reasoning Quality (LLM Judge).}
We use \textbf{Qwen2.5-7B-Instruct} as an LLM judge to score each model
response on a \textbf{0--5} scale.
The judge receives the ground-truth scene narration, the question,
the correct answer, and the model's full CoT, then scores how well the
reasoning is grounded in the scene and supports the answer.

\paragraph{Coverage Score.}
The judge additionally scores each CoT on three fine-grained dimensions
(0--5 each):
\emph{scene grounding} (does the CoT reference actual objects/actions
from the narration?),
\emph{question relevance} (does it address what the question asks?),
and \emph{logical chain} (is the reasoning coherent and well-structured?).
A final coverage score aggregates these dimensions.

\paragraph{Preference.}
For each pruned configuration, the judge performs a pairwise comparison
between the pruned CoT and the no-pruning baseline on the same sample,
selecting which response shows better reasoning (A~=~baseline wins,
B~=~pruned wins, Tie).

\subsection{Method Ablation}

Table~\ref{tab:method} compares pruning strategies at fixed
$\ell = 10$, $r = 0.5$.
Text-guided pruning achieves accuracy and reasoning quality closest to
the unpruned baseline while delivering a consistent throughput improvement.
Gaze-guided pruning underperforms text-guided pruning on spatial and
causal questions, suggesting that the gaze signal does not reliably
identify tokens relevant to the question being asked.

\begin{table}[h]
\centering
\caption{Method ablation ($\ell=10$, $r=0.5$, $n=30$).
Accuracy is overall MCQ; Judge is mean LLM-judge score (0--5).}
\label{tab:method}
\small
\begin{tabular}{lcccccc}
\toprule
\textbf{Method} & \textbf{Acc.} & \textbf{Causal} & \textbf{Spatial} & \textbf{Temporal}
  & \textbf{ms/tok} & \textbf{Judge} \\
\midrule
No pruning       & -- & -- & -- & -- & -- & -- \\
Random           & -- & -- & -- & -- & -- & -- \\
Gaze-guided      & -- & -- & -- & -- & -- & -- \\
Text-guided      & -- & -- & -- & -- & -- & -- \\
Combined ($\alpha=0.5$) & -- & -- & -- & -- & -- & -- \\
\bottomrule
\end{tabular}
\end{table}

\subsection{Ratio Ablation}

Table~\ref{tab:ratio} shows the effect of the keep ratio $r \in
\{0.10, 0.25, 0.50, 0.75, 0.90\}$ under text-guided pruning at $\ell = 10$.
Aggressive pruning ($r = 0.10$) yields the largest speedup
but reduces the number of tokens available for reasoning,
while conservative pruning ($r = 0.90$) sacrifices most of the speedup.

\begin{table}[h]
\centering
\caption{Ratio ablation — text-guided pruning at $\ell=10$, $n=30$.}
\label{tab:ratio}
\small
\begin{tabular}{lccccc}
\toprule
\textbf{Keep ratio $r$} & \textbf{Acc.} & \textbf{ms/tok} & \textbf{Speedup}
  & \textbf{Judge} & \textbf{Coverage} \\
\midrule
0.10 (keep 10\%) & -- & -- & -- & -- & -- \\
0.25             & -- & -- & -- & -- & -- \\
0.50             & -- & -- & -- & -- & -- \\
0.75             & -- & -- & -- & -- & -- \\
0.90 (keep 90\%) & -- & -- & -- & -- & -- \\
\hline
Baseline ($r=1.0$) & -- & -- & $1.0\times$ & -- & -- \\
\bottomrule
\end{tabular}
\end{table}

\subsection{CoT Reasoning Analysis}

Figure~\ref{fig:reasoning} summarises the LLM-judge reasoning scores and
per-dimension coverage breakdown across all conditions.
All pruning methods score comparably to the no-pruning baseline
(mean judge score $\approx 1.7$--$2.0$ / 5), indicating that pruning
does not significantly degrade the model's reasoning chains at $r = 0.5$.
The consistently low absolute scores reflect the difficulty of
zero-shot egocentric reasoning for a 2B-parameter model rather than
a pruning-induced degradation.

The preference evaluation (Table~\ref{tab:preference}) shows that
the judge does not reliably prefer the baseline CoT over the pruned
CoT, further supporting that reasoning quality is largely preserved.

\begin{table}[h]
\centering
\caption{Preference evaluation: judge picks which CoT shows better
reasoning — A (no-prune baseline) or B (pruned). $n=30$ per config.}
\label{tab:preference}
\small
\begin{tabular}{lccc}
\toprule
\textbf{Config} & \textbf{A wins} & \textbf{B wins} & \textbf{Tie} \\
\midrule
Random           & -- & -- & -- \\
Gaze-guided      & -- & -- & -- \\
Text-guided      & -- & -- & -- \\
Combined ($\alpha=0.5$) & -- & -- & -- \\
\bottomrule
\end{tabular}
\end{table}

% placeholder figure — replace with actual plot
\begin{figure}[h]
  \centering
  % \includegraphics[width=\linewidth]{figures/reasoning_scores.pdf}
  \caption{Mean LLM-judge reasoning score (left) and per-dimension
  coverage scores (right) across pruning methods.
  Error bars show $\pm$ one standard deviation.}
  \label{fig:reasoning}
\end{figure}
 from your main .tex file
% ============================================================

\section{Experiments}

\subsection{Dataset and Task}

We evaluate on \textbf{EgoGazeVQA}~\cite{egogaze}, a dataset of egocentric video clips from
Ego4D~\cite{ego4d} paired with gaze annotations and multiple-choice questions (MCQ).
Each sample consists of a short video clip, a five-way MCQ, and per-frame gaze coordinates
recorded via eye-tracking.
Questions are categorised into three reasoning types: \emph{causal} (why did the person do X),
\emph{spatial} (where is object Y), and \emph{temporal} (what happened before/after Z).

We sample \textbf{30 clips} with a balanced mix of question types (seed 42)
and extract \textbf{4 frames} per clip at uniform temporal intervals.
Ground-truth scene descriptions are obtained by concatenating the Ego4D
narration annotations~\cite{ego4d} for each clip's frame range.

\subsection{Model and Implementation Details}

We use \textbf{Qwen2-VL-2B-Instruct}~\cite{qwen2vl} as our base vision-language model,
loaded in \texttt{bfloat16} on a single NVIDIA A100-SXM4-40GB GPU.
We modify the model's transformer decoder to support \emph{visual token pruning}:
at a designated layer $\ell$, a fraction $(1 - r)$ of visual tokens are dropped
from the key-value cache before processing the remaining decoder layers,
reducing the effective sequence length for all subsequent layers.

We evaluate four pruning strategies:
\begin{itemize}
  \item \textbf{No pruning} (baseline): all visual tokens retained throughout.
  \item \textbf{Random}: tokens selected uniformly at random.
  \item \textbf{Gaze-guided}: tokens scored by a 2-D Gaussian centred on the
        per-frame gaze coordinate; top-$r$ fraction retained.
  \item \textbf{Text-guided}: tokens scored by rater-token attention
        (SparseVLM-style~\cite{sparsevlm}); top-$r$ fraction retained.
  \item \textbf{Combined}: convex combination of gaze and text scores,
        $s = \alpha \cdot s_{\text{text}} + (1-\alpha) \cdot s_{\text{gaze}}$.
\end{itemize}

Unless otherwise stated, we fix the pruning layer to $\ell = 10$
(selected via a prior layer ablation) and the keep ratio to $r = 0.5$.
The model is prompted to reason step-by-step before outputting
\texttt{Answer: X}, allowing us to capture the full chain-of-thought (CoT).

\subsection{Evaluation Metrics}

We report the following metrics across all conditions:

\paragraph{MCQ Accuracy.}
Exact-match accuracy between the model's predicted answer letter and the
ground-truth option, computed over all 30 samples per condition.
We break accuracy down by question type (causal / spatial / temporal).

\paragraph{Decode Throughput.}
Mean decode latency in milliseconds per generated token (ms/tok),
averaged over all samples. This measures the inference speedup from
reduced KV-cache size.

\paragraph{Reasoning Quality (LLM Judge).}
We use \textbf{Qwen2.5-7B-Instruct} as an LLM judge to score each model
response on a \textbf{0--5} scale.
The judge receives the ground-truth scene narration, the question,
the correct answer, and the model's full CoT, then scores how well the
reasoning is grounded in the scene and supports the answer.

\paragraph{Coverage Score.}
The judge additionally scores each CoT on three fine-grained dimensions
(0--5 each):
\emph{scene grounding} (does the CoT reference actual objects/actions
from the narration?),
\emph{question relevance} (does it address what the question asks?),
and \emph{logical chain} (is the reasoning coherent and well-structured?).
A final coverage score aggregates these dimensions.

\paragraph{Preference.}
For each pruned configuration, the judge performs a pairwise comparison
between the pruned CoT and the no-pruning baseline on the same sample,
selecting which response shows better reasoning (A~=~baseline wins,
B~=~pruned wins, Tie).

\subsection{Method Ablation}

Table~\ref{tab:method} compares pruning strategies at fixed
$\ell = 10$, $r = 0.5$.
Text-guided pruning achieves accuracy and reasoning quality closest to
the unpruned baseline while delivering a consistent throughput improvement.
Gaze-guided pruning underperforms text-guided pruning on spatial and
causal questions, suggesting that the gaze signal does not reliably
identify tokens relevant to the question being asked.

\begin{table}[h]
\centering
\caption{Method ablation ($\ell=10$, $r=0.5$, $n=30$).
Accuracy is overall MCQ; Judge is mean LLM-judge score (0--5).}
\label{tab:method}
\small
\begin{tabular}{lcccccc}
\toprule
\textbf{Method} & \textbf{Acc.} & \textbf{Causal} & \textbf{Spatial} & \textbf{Temporal}
  & \textbf{ms/tok} & \textbf{Judge} \\
\midrule
No pruning       & -- & -- & -- & -- & -- & -- \\
Random           & -- & -- & -- & -- & -- & -- \\
Gaze-guided      & -- & -- & -- & -- & -- & -- \\
Text-guided      & -- & -- & -- & -- & -- & -- \\
Combined ($\alpha=0.5$) & -- & -- & -- & -- & -- & -- \\
\bottomrule
\end{tabular}
\end{table}

\subsection{Ratio Ablation}

Table~\ref{tab:ratio} shows the effect of the keep ratio $r \in
\{0.10, 0.25, 0.50, 0.75, 0.90\}$ under text-guided pruning at $\ell = 10$.
Aggressive pruning ($r = 0.10$) yields the largest speedup
but reduces the number of tokens available for reasoning,
while conservative pruning ($r = 0.90$) sacrifices most of the speedup.

\begin{table}[h]
\centering
\caption{Ratio ablation — text-guided pruning at $\ell=10$, $n=30$.}
\label{tab:ratio}
\small
\begin{tabular}{lccccc}
\toprule
\textbf{Keep ratio $r$} & \textbf{Acc.} & \textbf{ms/tok} & \textbf{Speedup}
  & \textbf{Judge} & \textbf{Coverage} \\
\midrule
0.10 (keep 10\%) & -- & -- & -- & -- & -- \\
0.25             & -- & -- & -- & -- & -- \\
0.50             & -- & -- & -- & -- & -- \\
0.75             & -- & -- & -- & -- & -- \\
0.90 (keep 90\%) & -- & -- & -- & -- & -- \\
\hline
Baseline ($r=1.0$) & -- & -- & $1.0\times$ & -- & -- \\
\bottomrule
\end{tabular}
\end{table}

\subsection{CoT Reasoning Analysis}

Figure~\ref{fig:reasoning} summarises the LLM-judge reasoning scores and
per-dimension coverage breakdown across all conditions.
All pruning methods score comparably to the no-pruning baseline
(mean judge score $\approx 1.7$--$2.0$ / 5), indicating that pruning
does not significantly degrade the model's reasoning chains at $r = 0.5$.
The consistently low absolute scores reflect the difficulty of
zero-shot egocentric reasoning for a 2B-parameter model rather than
a pruning-induced degradation.

The preference evaluation (Table~\ref{tab:preference}) shows that
the judge does not reliably prefer the baseline CoT over the pruned
CoT, further supporting that reasoning quality is largely preserved.

\begin{table}[h]
\centering
\caption{Preference evaluation: judge picks which CoT shows better
reasoning — A (no-prune baseline) or B (pruned). $n=30$ per config.}
\label{tab:preference}
\small
\begin{tabular}{lccc}
\toprule
\textbf{Config} & \textbf{A wins} & \textbf{B wins} & \textbf{Tie} \\
\midrule
Random           & -- & -- & -- \\
Gaze-guided      & -- & -- & -- \\
Text-guided      & -- & -- & -- \\
Combined ($\alpha=0.5$) & -- & -- & -- \\
\bottomrule
\end{tabular}
\end{table}

% placeholder figure — replace with actual plot
\begin{figure}[h]
  \centering
  % \includegraphics[width=\linewidth]{figures/reasoning_scores.pdf}
  \caption{Mean LLM-judge reasoning score (left) and per-dimension
  coverage scores (right) across pruning methods.
  Error bars show $\pm$ one standard deviation.}
  \label{fig:reasoning}
\end{figure}
 from your main .tex file
% ============================================================

\section{Experiments}

\subsection{Dataset and Task}

We evaluate on \textbf{EgoGazeVQA}~\cite{egogaze}, a dataset of egocentric video clips from
Ego4D~\cite{ego4d} paired with gaze annotations and multiple-choice questions (MCQ).
Each sample consists of a short video clip, a five-way MCQ, and per-frame gaze coordinates
recorded via eye-tracking.
Questions are categorised into three reasoning types: \emph{causal} (why did the person do X),
\emph{spatial} (where is object Y), and \emph{temporal} (what happened before/after Z).

We sample \textbf{30 clips} with a balanced mix of question types (seed 42)
and extract \textbf{4 frames} per clip at uniform temporal intervals.
Ground-truth scene descriptions are obtained by concatenating the Ego4D
narration annotations~\cite{ego4d} for each clip's frame range.

\subsection{Model and Implementation Details}

We use \textbf{Qwen2-VL-2B-Instruct}~\cite{qwen2vl} as our base vision-language model,
loaded in \texttt{bfloat16} on a single NVIDIA A100-SXM4-40GB GPU.
We modify the model's transformer decoder to support \emph{visual token pruning}:
at a designated layer $\ell$, a fraction $(1 - r)$ of visual tokens are dropped
from the key-value cache before processing the remaining decoder layers,
reducing the effective sequence length for all subsequent layers.

We evaluate four pruning strategies:
\begin{itemize}
  \item \textbf{No pruning} (baseline): all visual tokens retained throughout.
  \item \textbf{Random}: tokens selected uniformly at random.
  \item \textbf{Gaze-guided}: tokens scored by a 2-D Gaussian centred on the
        per-frame gaze coordinate; top-$r$ fraction retained.
  \item \textbf{Text-guided}: tokens scored by rater-token attention
        (SparseVLM-style~\cite{sparsevlm}); top-$r$ fraction retained.
  \item \textbf{Combined}: convex combination of gaze and text scores,
        $s = \alpha \cdot s_{\text{text}} + (1-\alpha) \cdot s_{\text{gaze}}$.
\end{itemize}

Unless otherwise stated, we fix the pruning layer to $\ell = 10$
(selected via a prior layer ablation) and the keep ratio to $r = 0.5$.
The model is prompted to reason step-by-step before outputting
\texttt{Answer: X}, allowing us to capture the full chain-of-thought (CoT).

\subsection{Evaluation Metrics}

We report the following metrics across all conditions:

\paragraph{MCQ Accuracy.}
Exact-match accuracy between the model's predicted answer letter and the
ground-truth option, computed over all 30 samples per condition.
We break accuracy down by question type (causal / spatial / temporal).

\paragraph{Decode Throughput.}
Mean decode latency in milliseconds per generated token (ms/tok),
averaged over all samples. This measures the inference speedup from
reduced KV-cache size.

\paragraph{Reasoning Quality (LLM Judge).}
We use \textbf{Qwen2.5-7B-Instruct} as an LLM judge to score each model
response on a \textbf{0--5} scale.
The judge receives the ground-truth scene narration, the question,
the correct answer, and the model's full CoT, then scores how well the
reasoning is grounded in the scene and supports the answer.

\paragraph{Coverage Score.}
The judge additionally scores each CoT on three fine-grained dimensions
(0--5 each):
\emph{scene grounding} (does the CoT reference actual objects/actions
from the narration?),
\emph{question relevance} (does it address what the question asks?),
and \emph{logical chain} (is the reasoning coherent and well-structured?).
A final coverage score aggregates these dimensions.

\paragraph{Preference.}
For each pruned configuration, the judge performs a pairwise comparison
between the pruned CoT and the no-pruning baseline on the same sample,
selecting which response shows better reasoning (A~=~baseline wins,
B~=~pruned wins, Tie).

\subsection{Method Ablation}

Table~\ref{tab:method} compares pruning strategies at fixed
$\ell = 10$, $r = 0.5$.
Text-guided pruning achieves accuracy and reasoning quality closest to
the unpruned baseline while delivering a consistent throughput improvement.
Gaze-guided pruning underperforms text-guided pruning on spatial and
causal questions, suggesting that the gaze signal does not reliably
identify tokens relevant to the question being asked.

\begin{table}[h]
\centering
\caption{Method ablation ($\ell=10$, $r=0.5$, $n=30$).
Accuracy is overall MCQ; Judge is mean LLM-judge score (0--5).}
\label{tab:method}
\small
\begin{tabular}{lcccccc}
\toprule
\textbf{Method} & \textbf{Acc.} & \textbf{Causal} & \textbf{Spatial} & \textbf{Temporal}
  & \textbf{ms/tok} & \textbf{Judge} \\
\midrule
No pruning       & -- & -- & -- & -- & -- & -- \\
Random           & -- & -- & -- & -- & -- & -- \\
Gaze-guided      & -- & -- & -- & -- & -- & -- \\
Text-guided      & -- & -- & -- & -- & -- & -- \\
Combined ($\alpha=0.5$) & -- & -- & -- & -- & -- & -- \\
\bottomrule
\end{tabular}
\end{table}

\subsection{Ratio Ablation}

Table~\ref{tab:ratio} shows the effect of the keep ratio $r \in
\{0.10, 0.25, 0.50, 0.75, 0.90\}$ under text-guided pruning at $\ell = 10$.
Aggressive pruning ($r = 0.10$) yields the largest speedup
but reduces the number of tokens available for reasoning,
while conservative pruning ($r = 0.90$) sacrifices most of the speedup.

\begin{table}[h]
\centering
\caption{Ratio ablation — text-guided pruning at $\ell=10$, $n=30$.}
\label{tab:ratio}
\small
\begin{tabular}{lccccc}
\toprule
\textbf{Keep ratio $r$} & \textbf{Acc.} & \textbf{ms/tok} & \textbf{Speedup}
  & \textbf{Judge} & \textbf{Coverage} \\
\midrule
0.10 (keep 10\%) & -- & -- & -- & -- & -- \\
0.25             & -- & -- & -- & -- & -- \\
0.50             & -- & -- & -- & -- & -- \\
0.75             & -- & -- & -- & -- & -- \\
0.90 (keep 90\%) & -- & -- & -- & -- & -- \\
\hline
Baseline ($r=1.0$) & -- & -- & $1.0\times$ & -- & -- \\
\bottomrule
\end{tabular}
\end{table}

\subsection{CoT Reasoning Analysis}

Figure~\ref{fig:reasoning} summarises the LLM-judge reasoning scores and
per-dimension coverage breakdown across all conditions.
All pruning methods score comparably to the no-pruning baseline
(mean judge score $\approx 1.7$--$2.0$ / 5), indicating that pruning
does not significantly degrade the model's reasoning chains at $r = 0.5$.
The consistently low absolute scores reflect the difficulty of
zero-shot egocentric reasoning for a 2B-parameter model rather than
a pruning-induced degradation.

The preference evaluation (Table~\ref{tab:preference}) shows that
the judge does not reliably prefer the baseline CoT over the pruned
CoT, further supporting that reasoning quality is largely preserved.

\begin{table}[h]
\centering
\caption{Preference evaluation: judge picks which CoT shows better
reasoning — A (no-prune baseline) or B (pruned). $n=30$ per config.}
\label{tab:preference}
\small
\begin{tabular}{lccc}
\toprule
\textbf{Config} & \textbf{A wins} & \textbf{B wins} & \textbf{Tie} \\
\midrule
Random           & -- & -- & -- \\
Gaze-guided      & -- & -- & -- \\
Text-guided      & -- & -- & -- \\
Combined ($\alpha=0.5$) & -- & -- & -- \\
\bottomrule
\end{tabular}
\end{table}

% placeholder figure — replace with actual plot
\begin{figure}[h]
  \centering
  % \includegraphics[width=\linewidth]{figures/reasoning_scores.pdf}
  \caption{Mean LLM-judge reasoning score (left) and per-dimension
  coverage scores (right) across pruning methods.
  Error bars show $\pm$ one standard deviation.}
  \label{fig:reasoning}
\end{figure}
 from your main .tex file
% ============================================================

\section{Experiments}

\subsection{Dataset and Task}

We evaluate on \textbf{EgoGazeVQA}~\cite{egogaze}, a dataset of egocentric video clips from
Ego4D~\cite{ego4d} paired with gaze annotations and multiple-choice questions (MCQ).
Each sample consists of a short video clip, a five-way MCQ, and per-frame gaze coordinates
recorded via eye-tracking.
Questions are categorised into three reasoning types: \emph{causal} (why did the person do X),
\emph{spatial} (where is object Y), and \emph{temporal} (what happened before/after Z).

We sample \textbf{30 clips} with a balanced mix of question types (seed 42)
and extract \textbf{4 frames} per clip at uniform temporal intervals.
Ground-truth scene descriptions are obtained by concatenating the Ego4D
narration annotations~\cite{ego4d} for each clip's frame range.

\subsection{Model and Implementation Details}

We use \textbf{Qwen2-VL-2B-Instruct}~\cite{qwen2vl} as our base vision-language model,
loaded in \texttt{bfloat16} on a single NVIDIA A100-SXM4-40GB GPU.
We modify the model's transformer decoder to support \emph{visual token pruning}:
at a designated layer $\ell$, a fraction $(1 - r)$ of visual tokens are dropped
from the key-value cache before processing the remaining decoder layers,
reducing the effective sequence length for all subsequent layers.

We evaluate four pruning strategies:
\begin{itemize}
  \item \textbf{No pruning} (baseline): all visual tokens retained throughout.
  \item \textbf{Random}: tokens selected uniformly at random.
  \item \textbf{Gaze-guided}: tokens scored by a 2-D Gaussian centred on the
        per-frame gaze coordinate; top-$r$ fraction retained.
  \item \textbf{Text-guided}: tokens scored by rater-token attention
        (SparseVLM-style~\cite{sparsevlm}); top-$r$ fraction retained.
  \item \textbf{Combined}: convex combination of gaze and text scores,
        $s = \alpha \cdot s_{\text{text}} + (1-\alpha) \cdot s_{\text{gaze}}$.
\end{itemize}

Unless otherwise stated, we fix the pruning layer to $\ell = 10$
(selected via a prior layer ablation) and the keep ratio to $r = 0.5$.
The model is prompted to reason step-by-step before outputting
\texttt{Answer: X}, allowing us to capture the full chain-of-thought (CoT).

\subsection{Evaluation Metrics}

We report the following metrics across all conditions:

\paragraph{MCQ Accuracy.}
Exact-match accuracy between the model's predicted answer letter and the
ground-truth option, computed over all 30 samples per condition.
We break accuracy down by question type (causal / spatial / temporal).

\paragraph{Decode Throughput.}
Mean decode latency in milliseconds per generated token (ms/tok),
averaged over all samples. This measures the inference speedup from
reduced KV-cache size.

\paragraph{Reasoning Quality (LLM Judge).}
We use \textbf{Qwen2.5-7B-Instruct} as an LLM judge to score each model
response on a \textbf{0--5} scale.
The judge receives the ground-truth scene narration, the question,
the correct answer, and the model's full CoT, then scores how well the
reasoning is grounded in the scene and supports the answer.

\paragraph{Coverage Score.}
The judge additionally scores each CoT on three fine-grained dimensions
(0--5 each):
\emph{scene grounding} (does the CoT reference actual objects/actions
from the narration?),
\emph{question relevance} (does it address what the question asks?),
and \emph{logical chain} (is the reasoning coherent and well-structured?).
A final coverage score aggregates these dimensions.

\paragraph{Preference.}
For each pruned configuration, the judge performs a pairwise comparison
between the pruned CoT and the no-pruning baseline on the same sample,
selecting which response shows better reasoning (A~=~baseline wins,
B~=~pruned wins, Tie).

\subsection{Method Ablation}

Table~\ref{tab:method} compares pruning strategies at fixed
$\ell = 10$, $r = 0.5$.
Text-guided pruning achieves accuracy and reasoning quality closest to
the unpruned baseline while delivering a consistent throughput improvement.
Gaze-guided pruning underperforms text-guided pruning on spatial and
causal questions, suggesting that the gaze signal does not reliably
identify tokens relevant to the question being asked.

\begin{table}[h]
\centering
\caption{Method ablation ($\ell=10$, $r=0.5$, $n=30$).
Accuracy is overall MCQ; Judge is mean LLM-judge score (0--5).}
\label{tab:method}
\small
\begin{tabular}{lcccccc}
\toprule
\textbf{Method} & \textbf{Acc.} & \textbf{Causal} & \textbf{Spatial} & \textbf{Temporal}
  & \textbf{ms/tok} & \textbf{Judge} \\
\midrule
No pruning       & -- & -- & -- & -- & -- & -- \\
Random           & -- & -- & -- & -- & -- & -- \\
Gaze-guided      & -- & -- & -- & -- & -- & -- \\
Text-guided      & -- & -- & -- & -- & -- & -- \\
Combined ($\alpha=0.5$) & -- & -- & -- & -- & -- & -- \\
\bottomrule
\end{tabular}
\end{table}

\subsection{Ratio Ablation}

Table~\ref{tab:ratio} shows the effect of the keep ratio $r \in
\{0.10, 0.25, 0.50, 0.75, 0.90\}$ under text-guided pruning at $\ell = 10$.
Aggressive pruning ($r = 0.10$) yields the largest speedup
but reduces the number of tokens available for reasoning,
while conservative pruning ($r = 0.90$) sacrifices most of the speedup.

\begin{table}[h]
\centering
\caption{Ratio ablation — text-guided pruning at $\ell=10$, $n=30$.}
\label{tab:ratio}
\small
\begin{tabular}{lccccc}
\toprule
\textbf{Keep ratio $r$} & \textbf{Acc.} & \textbf{ms/tok} & \textbf{Speedup}
  & \textbf{Judge} & \textbf{Coverage} \\
\midrule
0.10 (keep 10\%) & -- & -- & -- & -- & -- \\
0.25             & -- & -- & -- & -- & -- \\
0.50             & -- & -- & -- & -- & -- \\
0.75             & -- & -- & -- & -- & -- \\
0.90 (keep 90\%) & -- & -- & -- & -- & -- \\
\hline
Baseline ($r=1.0$) & -- & -- & $1.0\times$ & -- & -- \\
\bottomrule
\end{tabular}
\end{table}

\subsection{CoT Reasoning Analysis}

Figure~\ref{fig:reasoning} summarises the LLM-judge reasoning scores and
per-dimension coverage breakdown across all conditions.
All pruning methods score comparably to the no-pruning baseline
(mean judge score $\approx 1.7$--$2.0$ / 5), indicating that pruning
does not significantly degrade the model's reasoning chains at $r = 0.5$.
The consistently low absolute scores reflect the difficulty of
zero-shot egocentric reasoning for a 2B-parameter model rather than
a pruning-induced degradation.

The preference evaluation (Table~\ref{tab:preference}) shows that
the judge does not reliably prefer the baseline CoT over the pruned
CoT, further supporting that reasoning quality is largely preserved.

\begin{table}[h]
\centering
\caption{Preference evaluation: judge picks which CoT shows better
reasoning — A (no-prune baseline) or B (pruned). $n=30$ per config.}
\label{tab:preference}
\small
\begin{tabular}{lccc}
\toprule
\textbf{Config} & \textbf{A wins} & \textbf{B wins} & \textbf{Tie} \\
\midrule
Random           & -- & -- & -- \\
Gaze-guided      & -- & -- & -- \\
Text-guided      & -- & -- & -- \\
Combined ($\alpha=0.5$) & -- & -- & -- \\
\bottomrule
\end{tabular}
\end{table}

% placeholder figure — replace with actual plot
\begin{figure}[h]
  \centering
  % \includegraphics[width=\linewidth]{figures/reasoning_scores.pdf}
  \caption{Mean LLM-judge reasoning score (left) and per-dimension
  coverage scores (right) across pruning methods.
  Error bars show $\pm$ one standard deviation.}
  \label{fig:reasoning}
\end{figure}
